% Hello Buddy -- Electronic Product Development Report
\documentclass[12pt,a4paper]{article}

\usepackage[margin=1in]{geometry}
\usepackage[T1]{fontenc}
\usepackage{mathptmx}            % Times New Roman font
\usepackage{array}
\usepackage{longtable}
\usepackage{enumitem}
\usepackage{tabularx}
\usepackage{ragged2e}
\usepackage{needspace}

\setlength{\parindent}{0pt}
\setlength{\parskip}{6pt}
\renewcommand{\arraystretch}{1.35}

% Column types
\newcolumntype{L}[1]{>{\RaggedRight\arraybackslash}p{#1}}
\newcolumntype{Y}{>{\RaggedRight\arraybackslash}X}

\begin{document}

% ====================================================================
%                           TITLE
% ====================================================================
\begin{center}
    {\Large\bfseries Electronic Product Development}\\[0.6cm]
    {\large\bfseries Product: Hello Buddy}\\[0.15cm]
    {\large Desktop Companion and Study Assistant Robot}\\[0.5cm]
    {\normalsize Group Number: 09 \qquad Lead: Ziaur Rahaman \qquad Date: August 17, 2026}
\end{center}

\vspace{0.3cm}
\hrule
\vspace{0.4cm}

% ====================================================================
%                             AIM
% ====================================================================
\section*{Aim}

To develop a comprehensive requirements specification for \textbf{Hello Buddy}, a stationary
desktop companion and study-assistant robot, by engaging multiple stakeholders using formal
techniques and documenting the requirements following ISO/IEC/IEEE 29148 and ISO 9001
guidelines.

% ====================================================================
%                    EQUIPMENT / SOFTWARE
% ====================================================================
\section*{Equipment / Software Required}

\begin{itemize}[leftmargin=*]
\item Laptop with Arduino IDE (2.x) and ESP32-S3 board support package installed.
\item ESP32-S3 N16R8 development board (16\,MB flash, 8\,MB PSRAM).
\item Hardware peripherals: GC9A01 round TFT display, SSD1306 OLED
      128$\times$64, DFPlayer Mini, INMP441 I2S microphone, VL53L0X ToF
      sensor, TTP223 capacitive touch sensor $\times$2, SG90 micro-servo
      $\times$4, LM386 audio amplifier, 8\,$\Omega$ speaker, 7.4\,V LiPo
      battery, UBEC 5\,V/3\,A regulator.
\item Cloud and API services: LLM inference via Groq, DeepSeek, Gemini API,
      or any locally hosted LLM (e.g., Ollama); Firebase Realtime Database
      for persistent memory and task storage; FastAPI backend (Python)
      running on a PC/server.
\item Android/iOS smartphone for the companion Flutter application.
\item 3D-printed enclosure (PLA filament, $\sim$150\,g) based on the
      EMOtto STL body.
\item Tools: soldering station, multimeter, USB-C cable, microSD card
      ($\leq$32\,GB, FAT32).
\item Reference: ISO/IEC/IEEE 29148:2018 for requirement attributes and
      documentation structure.
\end{itemize}

% ====================================================================
%                           THEORY
% ====================================================================
\section*{Theory}

After gathering requirements, the analyst reviews and analyses them to understand the
exact customer needs and resolve contradictions, ambiguities, and incomplete information.

This is necessary because stakeholders often have partial views of the system.
The analyst engages in further discussions to address these problems.

Effective requirements analysis involves answering key questions, such as:

\begin{itemize}[leftmargin=*]
\item What problem does Hello Buddy solve, and why is physical presence
      important?
\item What are the system inputs (voice, touch, proximity) and outputs
      (expressions, audio, OLED, servo motion)?
\item What procedures are needed for the Pomodoro study loop and task
      management?
\item What complexities arise from simultaneous display rendering, audio
      playback, and API calls?
\item How will Hello Buddy interface with external LLM APIs (Groq,
      DeepSeek, Gemini, or a local LLM), Firebase, and the Flutter app?
\end{itemize}

Once the customer requirements are clearly understood, the analyst identifies and
resolves anomalies, including ambiguities that could cause misinterpretations during
development.

% ====================================================================
%                          PROCEDURE
% ====================================================================
\newpage
\section*{Procedure}

\subsection*{Pre-Requirements Activities}

\textbf{a.\ Stakeholder Identification}

\begin{enumerate}[leftmargin=*]
\item List all stakeholders: students (primary users), Faculty-in-Charge,
      firmware engineers, Flutter app developers, hardware assembly
      engineers, and regulatory/safety reviewers.
\item Use organisational charts, brainstorming sessions, and focus-group
      discussions.
\end{enumerate}

\vspace{0.3cm}

\textbf{Stakeholder Examples}

\begin{center}
\begin{tabular}{|L{0.28\textwidth}|L{0.58\textwidth}|}
\hline
\textbf{Stakeholder Type} & \textbf{Examples} \\ \hline
Customer / User & Students using Hello Buddy while studying
                  (interviewed: Atanu Hazra, B.Tech 2nd Year;
                  Arka Ghosh, B.Tech 3rd Year) \\ \hline
Technical Staff & Firmware engineers (Arduino/FreeRTOS), app
                  developers (Flutter), PCB designers \\ \hline
Management & Project lead, QA reviewer, Faculty-in-Charge \\ \hline
Regulatory & Compliance officers reviewing CE/FCC,
             IEC 62133 (LiPo battery safety) \\ \hline
\end{tabular}
\end{center}

\subsection*{Requirements Elicitation Techniques}

\textbf{Structured Interviews}
\begin{itemize}[leftmargin=*]
\item Two structured user-validation interviews were conducted with
      students across different academic years to understand real desk
      habits, pain points, and expectations.
\item Each interview comprised 15 questions covering study behaviour,
      product reaction, feature preferences, design opinions, and
      willingness to pay.
\item Responses were documented verbatim and analysed for recurring
      patterns before the group discussion.
\end{itemize}

\textbf{Group Discussions / Workshops}
\begin{itemize}[leftmargin=*]
\item A formal group discussion was held after the interviews, using the
      interview findings as the primary reference.
\item All four members participated: Ziaur Rahaman, Aryan Raj, Shamil
      Minhaj, and Ayush Dutta.
\item Key design decisions (stationary design, customisable expressions,
      wake-word microphone, Pomodoro integration) were formally agreed
      upon and documented in the group discussion minutes.
\end{itemize}

\textbf{Questionnaires / Surveys}
\begin{itemize}[leftmargin=*]
\item Used for broader stakeholder input to gauge preferred price range,
      feature priority, and willingness to use voice commands.
\item Included both qualitative and quantitative questions.
\end{itemize}

\textbf{Observation and Site Visits}
\begin{itemize}[leftmargin=*]
\item Observed real-world desk study behaviour and distraction patterns
      among students.
\item Identified that phone distraction (social media and short-form
      video) was the single biggest obstacle to sustained focus, as
      confirmed by both interviewees.
\end{itemize}

\textbf{Review Existing Documentation}
\begin{itemize}[leftmargin=*]
\item Studied similar open-source projects (Otto DIY, EMOtto\_ESP32 on
      Printables.com) and competitor products (Loona, Vector, ElliQ)
      for feature benchmarking.
\end{itemize}

% ====================================================================
%              KEY FINDINGS FROM USER INTERVIEWS
% ====================================================================
\newpage
\subsection*{Key Findings from User Interviews}

The following consistent patterns emerged across both interviews
(Atanu Hazra, B.Tech 2nd Year; Arka Ghosh, B.Tech 3rd Year):

\begin{itemize}[leftmargin=*]
\item Both participants spent 4 to 7 hours at their desks daily and
      described their workspaces as monotonous and uninspiring.
\item Phone distraction was identified as the single biggest obstacle
      to sustained focus, with social media and short-form video being
      the primary culprits.
\item Both responded positively to a physical, expressive companion on
      the desk, preferring it over another app-based solution.
\item A round, face-like display was strongly preferred over a flat
      rectangular screen, as it communicated personality rather than
      functionality alone.
\item Customisable expressions and sounds were considered essential to
      prevent the product from feeling generic over time.
\item Both raised concerns about the microphone being always active,
      requesting a wake-word system with a visible LED indicator and
      a physical disable button.
\item A stable, stationary placement was preferred over a walking or
      movable design, as both cited the risk of the robot falling from
      the desk and damaging itself.
\item Reward-based expressions (e.g., a proud or happy face after
      completing a study session) were identified as a strong
      motivational feature.
\item Perceived value ranged from Rs.\,3{,}500 to Rs.\,4{,}500, which
      is above the planned retail price of Rs.\,3{,}000--- confirming
      a strong value proposition.
\end{itemize}

These findings directly shaped the group discussion decisions and all
subsequent design requirements documented in this report.

% ====================================================================
%                    SAMPLE INTERVIEW TEMPLATE
% ====================================================================
\subsection*{Sample Interview Template (Based on Conducted Interviews)}

\begin{center}
\scriptsize
\begin{tabular}{|L{0.28\textwidth}|L{0.14\textwidth}|L{0.24\textwidth}|L{0.24\textwidth}|}
\hline
\textbf{Question} & \textbf{Stakeholder} &
\textbf{Response (Atanu Hazra)} &
\textbf{Response (Arka Ghosh)} \\ \hline
How many hours do you spend at your desk daily? &
User (Student) &
4 to 6 hours; more during exams. &
5 to 7 hours; up to 10 during deadlines. \\ \hline
Does your desk feel motivating or boring? &
User (Student) &
Boring mostly. Just a laptop and books. &
Mostly boring. Feels more like a workstation
than a place I want to be at. \\ \hline
Do you get distracted easily while studying? &
User (Student) &
Yes, constantly. Phone is the biggest problem. &
Yes. Instagram and YouTube--- open for one minute,
lose 30 minutes. \\ \hline
What is your first reaction to Hello Buddy? &
User (Student) &
Sounds really cool. The auto wake-up is a
nice detail. &
Really interesting. Something physical on the
desk feels different from just another app. \\ \hline
What feature excites you the most? &
User (Student) &
Face expressions and interactivity. Makes it
feel alive. &
Study timer with expressions. A happy face after
completing a session is motivating. \\ \hline
Would you prefer a round face or a flat screen? &
User (Student) &
Round. Feels more like a face. Rectangle
feels like a device. &
Round, for sure. Reads as a character rather
than a gadget. \\ \hline
Movable legs or stable single-place design? &
User (Student) &
Stable. If it falls from the table it could
get damaged. &
Stable. Movable parts mean more things that can
break. Risky near a laptop. \\ \hline
Does anything make you uncomfortable? &
User (Student) &
The microphone. Don't want it always listening. &
Always-on mic is a concern. Want a physical LED
indicator and a disable button. \\ \hline
What would make you stop using it after a week? &
User (Student) &
Bad battery life and repetitive sounds I cannot
change. &
Repetitive and predictable responses. Also poor
battery backup. \\ \hline
How much would you pay? &
User (Student) &
Around Rs.\,4{,}000 or more if it actually works
properly. &
Rs.\,3{,}500 to 4{,}500 depending on build quality
and how well features work. \\ \hline
\end{tabular}
\end{center}

\vspace{0.3cm}

\subsection*{Group Discussion Decisions}

Based on the interview insights, the group formally agreed on the
following design decisions:

\begin{enumerate}[leftmargin=*]
\item Hello Buddy will be a \textbf{stationary desktop device} with a
      round, face-like display and an optional subtle tilt mechanism
      for physical expressiveness.
\item The facial expression system will support both \textbf{graphical
      and ASCII-style} displays (e.g., \texttt{\^{}\_\^{}},
      \texttt{-\_-}, \texttt{T\_T}), with user-customisable expression
      sets and sound profiles.
\item A \textbf{Pomodoro-integrated study timer} will be a core
      built-in feature, with emotional expression feedback tied to
      each session phase (work, short break, long break).
\item The voice system will use a \textbf{wake-word model} with a
      hardware LED indicator and a physical microphone disable button
      to address privacy concerns.
\item A \textbf{companion mobile application} will manage
      customisation, task tracking, reminders, and session analytics.
\item \textbf{Power management} will be prioritised with an automatic
      sleep mode when the desk is unoccupied, detected via the
      VL53L0X proximity sensor.
\end{enumerate}

% ====================================================================
%                 CLASSIFYING REQUIREMENTS
% ====================================================================
\section*{Classifying Requirements Using Taxonomy}

Requirements are organised into a taxonomy to avoid duplication and ensure full
coverage. For Hello Buddy, requirements are classified into hardware, firmware,
cloud/API, and application layers.

\subsection*{Requirements Taxonomy Table}

\begin{center}
\small
\begin{tabular}{|L{0.13\textwidth}|L{0.15\textwidth}|L{0.38\textwidth}|L{0.20\textwidth}|}
\hline
\textbf{Category} & \textbf{Sub-Category} &
\textbf{Examples for Hello Buddy} &
\textbf{Type} \\ \hline
Hardware & Display &
GC9A01 round TFT must render animated expressions at
$\geq$10\,FPS. &
Functional \\ \hline
Hardware & Sensors &
VL53L0X must detect user presence within 10--60\,cm. &
Functional \\ \hline
Hardware & Power &
Device must operate $\geq$6\,h on 7.4\,V 1200\,mAh LiPo. &
Non-Functional (Reliability) \\ \hline
Hardware & Audio &
DFPlayer Mini must play MP3 files with $<$300\,ms latency. &
Non-Functional (Performance) \\ \hline
Firmware & Task Scheduling &
FreeRTOS must manage display, audio, servo, and WiFi
concurrently without deadlock. &
Non-Functional (Stability) \\ \hline
Firmware & Wake Word &
System must detect ``Hey Buddy'' within 2\,s of utterance. &
Functional \\ \hline
Cloud/API & LLM Integration &
LLM API (Groq/DeepSeek/Gemini/local) must return a JSON
action plan within 2\,s over WiFi. &
Non-Functional (Performance) \\ \hline
Cloud/API & Persistence &
Firebase Realtime DB must sync task list with
$<$500\,ms latency. &
Non-Functional (Performance) \\ \hline
Application & UI &
Flutter app must display live task list, Pomodoro timer,
and Hello Buddy status on one dashboard. &
Functional \\ \hline
Application & Security &
All API keys must be stored server-side; never exposed
to the Flutter app. &
Non-Functional (Security) \\ \hline
\end{tabular}
\end{center}

% ====================================================================
%              REQUIREMENTS DOCUMENTATION
% ====================================================================
\newpage
\section*{Requirements Documentation}

\subsection*{Functional vs.\ Non-Functional Requirements Taxonomy}

\begin{center}
\begin{tabular}{|L{0.25\textwidth}|L{0.63\textwidth}|}
\hline
\textbf{Type} & \textbf{Description / Examples} \\ \hline
Functional & Wake-word detection, proximity auto-wake, touch reaction,
             face expression change, Pomodoro timer, voice task addition,
             Spotify playback trigger, OLED status display \\ \hline
Non-Functional & API response $<$2\,s; battery $\geq$6\,h; concurrent
                 FreeRTOS tasks without jitter; no continuous microphone
                 recording \\ \hline
External Interface & LLM API (Groq/DeepSeek/Gemini/local via REST/HTTPS),
                     Firebase Realtime DB (WebSocket), Spotify Web API,
                     Flutter app over local HTTP/WebSocket \\ \hline
Design Constraints & ESP32-S3 Arduino IDE/FreeRTOS; fits inside EMOtto STL
                     body; 3.3\,V/5\,V circuits only; stationary desktop
                     design \\ \hline
\end{tabular}
\end{center}

\subsection*{Capturing Requirements Attributes (ISO/IEC/IEEE 29148)}

Each requirement carries the following attributes:

\begin{center}
\begin{tabular}{|L{0.25\textwidth}|L{0.63\textwidth}|}
\hline
\textbf{Attribute} & \textbf{Description} \\ \hline
Id & Unique identifier (e.g., HW-001, FW-002, API-001, APP-001) \\ \hline
Heading & Short name (e.g., ``GC9A01 Frame Rate'') \\ \hline
Text & Intent / description of the requirement \\ \hline
Owner & Who maintains this requirement \\ \hline
Priority & Must / Should / Could / Won't (MoSCoW) \\ \hline
Source & User interview, group discussion, or regulatory review \\ \hline
Rationale & Reason for inclusion \\ \hline
Difficulty & Estimated implementation difficulty \\ \hline
Type & Functional / Non-Functional / Interface / Constraint \\ \hline
Status & Draft / Reviewed / Approved \\ \hline
Verification Method & Test, Demonstration, Inspection, Analysis \\ \hline
\end{tabular}
\end{center}

% ====================================================================
%                    AFFINITY DIAGRAM
% ====================================================================
\section*{Affinity Diagram}

Brainstorming sessions were conducted with all four group members to gather diverse
input. The stated problem was: \textit{``How can a physical desk object help a student
stay focused, accountable, and motivated without becoming a distraction itself?''}

Participants wrote ideas on digital sticky notes, presented them, and grouped them on
a shared board. After multiple rounds, ideas were organised into five clusters:

\begin{enumerate}[leftmargin=*]
\item \textbf{Physical Presence and Accountability} --- proximity detection,
      worried face when user is inactive, passive monitoring via VL53L0X.
\item \textbf{Emotional Expression} --- animated face on GC9A01, customisable
      ASCII expressions (e.g., \texttt{\^{}\_\^{}}, \texttt{-\_-},
      \texttt{T\_T}), state-based emotion mapping.
\item \textbf{Study Assistance} --- Pomodoro timer, voice task management,
      OLED task display, Spotify playlist control.
\item \textbf{Connectivity} --- LLM API (Groq/DeepSeek/Gemini/local) for
      intelligence, Firebase for memory, Flutter app for control and
      statistics.
\item \textbf{Safety and Privacy} --- wake-word-only microphone activation,
      no continuous recording, voltage-limited circuits.
\end{enumerate}

Using affinity diagrams, Voices of Customer (VOCs) from user interviews were mapped
to higher-level needs, which fed into the Kano model classification below.

% ====================================================================
%                           AHP
% ====================================================================
\newpage
\section*{Analytical Hierarchy Process (AHP)}

Saaty's AHP was applied to rank the top customer needs identified during elicitation.
Pairwise comparisons were made between five high-level needs using the standard
1--9 scale.

\textbf{High-level needs compared:}
\begin{enumerate}[leftmargin=*]
\item Expressive face display ($N_1$)
\item Study assistance and task management ($N_2$)
\item Physical accountability and presence ($N_3$)
\item Voice interaction ($N_4$)
\item App connectivity and statistics ($N_5$)
\end{enumerate}

\textbf{Summary of AHP weightings (derived from group evaluation):}

\begin{center}
\begin{tabular}{|L{0.45\textwidth}|L{0.20\textwidth}|}
\hline
\textbf{Customer Need} & \textbf{Relative Weight} \\ \hline
Expressive face display & 0.30 \\ \hline
Physical accountability and presence & 0.25 \\ \hline
Study assistance and task management & 0.22 \\ \hline
App connectivity and statistics & 0.14 \\ \hline
Voice interaction & 0.09 \\ \hline
\end{tabular}
\end{center}

The expressive face display received the highest weight, consistent with user interview
feedback: ``if the face looks wrong or dead, the whole product feels wrong.''
Customer surveys and evaluations fed into the process, with care taken to avoid bias
or misleading data.

% ====================================================================
%                        KANO MODEL
% ====================================================================
\section*{Kano Model Integration}

After ranking, needs are classified into three categories:

\begin{itemize}[leftmargin=*]
\item \textbf{Exciters (Delighters):} Servo-based body reaction on touch;
      real-time emotion detection from voice tone via LLM API; customisable
      ASCII facial expressions; Spotify playlist auto-play when a study
      session starts.
\item \textbf{Expected Needs:} Wake-word detection before microphone activates;
      proximity-based auto-wake; DFPlayer audio reactions to events;
      OLED task display.
\item \textbf{Normal (Performance) Needs:} Battery life (more hours is better);
      API response latency (lower is better); number of face expressions
      available (more variety reduces repetitive behaviour).
\end{itemize}

% ====================================================================
%             DOCUMENTING REQUIREMENTS PER ISO
% ====================================================================
\section*{Documenting Hardware and Software Requirements per ISO Guidelines}

\begin{itemize}[leftmargin=*]
\item Follow ISO/IEC/IEEE 29148:2018 for documentation.
\item Requirements must be: Necessary, unambiguous, consistent, complete,
      singular (atomic), verifiable, feasible, traceable.
\item Structure: Use a hierarchical format with unique IDs
      (HW-, FW-, API-, APP-).
\item Validation: All requirements reviewed and approved by group members
      and Faculty-in-Charge.
\end{itemize}

% ====================================================================
%                HARDWARE REQUIREMENTS
% ====================================================================
\newpage
\section*{Hardware Requirements Documentation}

Hardware specifications focus on the ESP32-S3 N16R8 board and all connected
peripherals.

\begin{enumerate}[leftmargin=*]
\item \textbf{Gather and Refine:} Hardware needs identified via structured
      interviews with team engineers and user feedback sessions.
\item \textbf{Document in Template:} Use a table format with ISO attributes.
\end{enumerate}

\subsection*{Hardware Requirements Template Table}

\begin{center}
\scriptsize
\begin{tabular}{|L{0.07\textwidth}|L{0.35\textwidth}|L{0.10\textwidth}|L{0.18\textwidth}|L{0.18\textwidth}|}
\hline
\textbf{Req.~ID} & \textbf{Description} & \textbf{Priority} &
\textbf{Verification} & \textbf{Traceability} \\ \hline
HW-001 &
GC9A01 round TFT display must render animated face expressions at
$\geq$10\,FPS via SPI on GPIO 5/17/16/18/23. &
Must &
Performance Testing &
User Interview \\ \hline
HW-002 &
SSD1306 OLED 128$\times$64 must display task name, session time,
tasks remaining, and Pomodoro count via I2C on GPIO 39/40. &
Must &
Functional Testing &
Group Discussion \\ \hline
HW-003 &
INMP441 I2S microphone must capture audio on GPIO 41/42/2;
active only after wake-word trigger, not continuously. &
Must &
Inspection / Privacy Audit &
User Interview \\ \hline
HW-004 &
VL53L0X ToF sensor must detect user presence in 10--60\,cm
range via I2C on GPIO 22/21 with $\leq$50\,ms refresh. &
Must &
Lab Testing &
Group Discussion \\ \hline
HW-005 &
DFPlayer Mini must play MP3 responses from microSD within
$<$300\,ms of trigger via UART on GPIO 32 (RX, 1\,k$\Omega$)
/ GPIO 4 (TX). &
Must &
Timing Test &
User Interview \\ \hline
HW-006 &
Four SG90 servos on hardware PWM pins must execute body-tilt
and reaction routines without jitter under simultaneous audio
playback. &
Must &
Stress Testing &
Group Discussion \\ \hline
HW-007 &
7.4\,V 1200\,mAh LiPo through UBEC 5\,V/3\,A must sustain
full operation for $\geq$6 hours. &
Should &
Duration Test &
User Interview \\ \hline
HW-008 &
All circuits must operate at 3.3\,V or 5\,V only; LiPo must
connect via over-current protection (UBEC). &
Must &
Inspection &
Regulatory Review \\ \hline
\end{tabular}
\end{center}

% ====================================================================
%          FIRMWARE REQUIREMENTS
% ====================================================================
\newpage
\section*{Software / Firmware / API Requirements Documentation}

Software specifications cover the ESP32-S3 firmware, the FastAPI backend server,
and the Flutter companion application. All natural language processing is handled
by an \textbf{LLM API} (Groq, DeepSeek, Gemini, or any locally hosted LLM such
as Ollama) over WiFi; no on-device language model is used.

\begin{enumerate}[leftmargin=*]
\item \textbf{Gather and Refine:} Requirements elicited through group
      discussions with developers and feedback from user interviews.
\item \textbf{Document in Template:} Similar to hardware, but include API
      endpoints and use cases.
\end{enumerate}

\subsection*{Firmware Requirements (ESP32-S3)}

\begin{center}
\scriptsize
\begin{tabular}{|L{0.07\textwidth}|L{0.35\textwidth}|L{0.10\textwidth}|L{0.18\textwidth}|L{0.18\textwidth}|}
\hline
\textbf{Req.~ID} & \textbf{Description} & \textbf{Priority} &
\textbf{Verification} & \textbf{Traceability} \\ \hline
FW-001 &
Firmware must use FreeRTOS to run display, audio, servo, sensor
polling, and WiFi tasks concurrently without priority inversion
or deadlock. &
Must &
Stress Testing &
Group Discussion \\ \hline
FW-002 &
MicroWakeWord library must detect ``Hey Buddy'' and activate the
INMP441 microphone; microphone must remain inactive at all other
times. &
Must &
Functional Testing &
User Interview \\ \hline
FW-003 &
On wake-word detection, firmware must record audio, send to
FastAPI for Whisper STT, receive LLM API JSON action plan, and
execute within 3\,s end-to-end. &
Must &
Timing Test &
Group Discussion \\ \hline
FW-004 &
GC9A01 display must map at least 8 emotion states (Happy, Sad,
Focused, Stressed, Sleepy, Excited, Neutral, Task Complete) to
animated sprite sequences. &
Must &
Demonstration &
User Interview \\ \hline
FW-005 &
OLED must display scrolling task name, session countdown, tasks
remaining, and Pomodoro count; updated every second during an
active session. &
Must &
Functional Testing &
Group Discussion \\ \hline
FW-006 &
Touch on TTP223 (GPIO 12/33) must trigger face reaction and
DFPlayer audio response within $<$200\,ms. &
Must &
Timing Test &
User Interview \\ \hline
FW-007 &
VL53L0X must poll every 50\,ms; if no presence detected for
$>$2\,min during an active session, Hello Buddy must show
``worried'' face and play an audio check-in. &
Should &
Functional Testing &
Group Discussion \\ \hline
FW-008 &
Firmware must expose an HTTP/WebSocket server on port 80 for
direct commands: \texttt{POST /emotion}, \texttt{POST /oled},
\texttt{POST /move}, \texttt{GET /status}. &
Must &
Integration Testing &
Group Discussion \\ \hline
\end{tabular}
\end{center}

% ====================================================================
%            BACKEND API REQUIREMENTS
% ====================================================================
\newpage
\subsection*{Backend API Requirements (FastAPI on PC / Server)}

\begin{center}
\scriptsize
\begin{tabular}{|L{0.07\textwidth}|L{0.35\textwidth}|L{0.10\textwidth}|L{0.18\textwidth}|L{0.18\textwidth}|}
\hline
\textbf{Req.~ID} & \textbf{Description} & \textbf{Priority} &
\textbf{Verification} & \textbf{Traceability} \\ \hline
API-001 &
FastAPI backend must receive audio from ESP32-S3, run Whisper
STT locally, and return a text transcript within 1.5\,s. &
Must &
Timing Test &
Group Discussion \\ \hline
API-002 &
Backend must forward the transcript to the configured LLM API
(Groq, DeepSeek, Gemini, or local Ollama), parse the JSON
action plan, and relay commands to the ESP32-S3. &
Must &
Integration Testing &
Group Discussion \\ \hline
API-003 &
All tasks must be stored and synced using Firebase Realtime
Database; the Flutter app, ESP32-S3, and FastAPI backend must
read/write from the same Firebase project. &
Must &
Integration Testing &
Group Discussion \\ \hline
API-004 &
Backend must expose task CRUD endpoints:
\texttt{POST /tasks}, \texttt{GET /tasks},
\texttt{PATCH /tasks/\{id\}}, \texttt{DELETE /tasks/\{id\}}. &
Must &
Functional Testing &
Group Discussion \\ \hline
API-005 &
Backend must expose \texttt{POST /session/start} and
\texttt{POST /session/end} to log study sessions with subject,
duration, and Pomodoro count to Firebase. &
Must &
Functional Testing &
Group Discussion \\ \hline
API-006 &
LLM API keys and Spotify credentials must be stored only in
server environment variables; never transmitted to the Flutter
app or ESP32-S3. &
Must &
Code Review / Audit &
Regulatory Review \\ \hline
API-007 &
Backend must support Spotify control via Spotipy:
\texttt{POST /spotify/play}, \texttt{POST /spotify/pause},
and playlist selection by session state. &
Could &
Functional Testing &
Group Discussion \\ \hline
\end{tabular}
\end{center}

\vspace{0.5cm}

\subsection*{Flutter Application Requirements}

\begin{center}
\scriptsize
\begin{tabular}{|L{0.07\textwidth}|L{0.35\textwidth}|L{0.10\textwidth}|L{0.18\textwidth}|L{0.18\textwidth}|}
\hline
\textbf{Req.~ID} & \textbf{Description} & \textbf{Priority} &
\textbf{Verification} & \textbf{Traceability} \\ \hline
APP-001 &
App must display a home dashboard with live session timer,
current task, Pomodoro count, and Hello Buddy connection
status. &
Must &
UI Testing &
Group Discussion \\ \hline
APP-002 &
App must allow the user to add, edit, delete, and mark tasks
complete; changes must sync to Firebase within 1\,s. &
Must &
Functional Testing &
User Interview \\ \hline
APP-003 &
App must allow the user to start and stop a Pomodoro study
session (25\,min study / 5\,min break) and send session state
to the FastAPI backend. &
Must &
Functional Testing &
Group Discussion \\ \hline
APP-004 &
App must provide a Hello Buddy Control panel for manually
changing face expression, triggering a movement, and viewing
device status. &
Should &
UI Demonstration &
Group Discussion \\ \hline
APP-005 &
App must display a statistics screen: hours studied per day,
task completion rate, mood log over time, and a streak
calendar; data sourced from Firebase. &
Should &
Functional Testing &
Group Discussion \\ \hline
\end{tabular}
\end{center}

\vspace{0.3cm}

For integrated systems, cross-references are maintained (e.g., FW-003 traces to
API-002 for LLM routing; API-003 traces to APP-002 for Firebase task sync).

% ====================================================================
%            REQUIREMENTS ANALYSIS & VERIFICATION
% ====================================================================
\newpage
\section*{Requirements Analysis \& Verification}

\textbf{a.\ Ensure requirements are:}
\begin{enumerate}[leftmargin=*]
\item Clear, unambiguous, and complete.
\item Testable (suggest acceptance criteria, e.g., ``LLM API response
      $<$2\,s'', ``touch reaction $<$200\,ms'').
\item Traceable (unique IDs, cross-reference design/test phases).
\end{enumerate}

\textbf{b.}\ Conduct peer reviews with group members and Faculty-in-Charge
to validate correctness and completeness.

\textbf{c.}\ Resolve conflicts (e.g., servo jitter vs.\ audio playback)
via FreeRTOS task priority assignment and hardware PWM pin allocation.

% ====================================================================
%           ISO-COMPLIANT DOCUMENTATION STRUCTURE
% ====================================================================
\section*{ISO-Compliant Documentation Structure}

Following ISO/IEC/IEEE 29148 template sections:

\begin{enumerate}[leftmargin=*]
\item \textbf{Introduction}
    \begin{itemize}
    \item Purpose: Define requirements for Hello Buddy for development
          and validation.
    \item Scope: ESP32-S3 firmware, FastAPI backend, Firebase integration,
          Flutter app.
    \item Definitions: FreeRTOS, LLM API (Groq/DeepSeek/Gemini/local),
          VL53L0X, GC9A01, DFPlayer Mini, Pomodoro.
    \end{itemize}

\item \textbf{Overall Description}
    \begin{itemize}
    \item Product Perspective: A stationary desktop companion combining
          physical emotional expression with cloud-backed study assistance.
    \item User Needs: Accountability, focus support, task management,
          expressive interaction, and voice-activated convenience without
          interrupting study flow.
    \end{itemize}

\item \textbf{Specific Requirements}
    \begin{itemize}
    \item Functional Requirements: Wake-word, proximity detection, face
          expressions, task CRUD, Pomodoro sessions, Spotify control,
          OLED display.
    \item Non-Functional Requirements: Latency, battery life, concurrent
          task stability, privacy-by-design microphone activation.
    \item External Interface Requirements: LLM API, Firebase Realtime DB,
          Spotify Web API, Flutter app over HTTP/WebSocket.
    \item Design Constraints: Arduino IDE/FreeRTOS; 3.3\,V/5\,V only;
          fits EMOtto STL body; stationary desktop design; no mains power.
    \item Hardware Requirements: HW-001 through HW-008.
    \item Software/Firmware/API Requirements: FW-001--FW-008,
          API-001--API-007, APP-001--APP-005.
    \end{itemize}

\item \textbf{Control \& Retention (ISO 9001)}
    \begin{itemize}
    \item Maintain documented information: records, change history,
          rationale.
    \item Control access and versioning via shared repository.
    \item Retain required documentation for audits and compliance.
    \end{itemize}

\item \textbf{Review \& Sign-off}
    \begin{itemize}
    \item Requirements specification reviewed by all key stakeholders.
    \item Record sign-off with date and revision number.
    \end{itemize}
\end{enumerate}

% ====================================================================
%              SUMMARY TAXONOMY TABLE
% ====================================================================
\newpage
\section*{Sample Table: Hardware and Software Requirements Taxonomy}

\begin{center}
\begin{tabular}{|L{0.23\textwidth}|L{0.30\textwidth}|L{0.37\textwidth}|}
\hline
\textbf{Category} & \textbf{Subcategories} &
\textbf{Example Requirements} \\ \hline
Hardware Requirements &
Display, Sensors, Power, Audio, Motion &
GC9A01 $\geq$10\,FPS; VL53L0X 10--60\,cm; LiPo $\geq$6\,h;
DFPlayer $<$300\,ms; SG90 jitter-free PWM \\ \hline
Firmware Requirements &
FreeRTOS Tasks, Wake Word, GPIO &
Concurrent task scheduling; ``Hey Buddy'' detection;
touch $<$200\,ms response \\ \hline
API / Cloud Requirements &
LLM API, Firebase, Spotify, Whisper &
LLM response $<$2\,s; Firebase sync $<$500\,ms;
Spotify playlist control; STT $<$1.5\,s \\ \hline
Application Requirements &
Dashboard, Tasks, Stats, Control Panel &
Real-time task sync; Pomodoro control; mood log;
Hello Buddy manual control \\ \hline
\end{tabular}
\end{center}

% ====================================================================
%                       EVALUATION
% ====================================================================
\section*{Evaluation}

\begin{center}
\begin{tabular}{|L{0.35\textwidth}|L{0.45\textwidth}|}
\hline
\textbf{Group} & \\ \hline
Member 1 name & Ziaur Rahaman \\ \hline
Member 2 name & Aryan Raj \\ \hline
Member 3 name & Ayush Dutta \\ \hline
Member 4 name & Shamil Minhaj \\ \hline
\end{tabular}
\end{center}

\textbf{Note:} Submit the report to FIC and upload the PDF copy of the same to
Google Classroom after evaluation.

\end{document}
